\scanpage{113}
\textcircled{2'}\quad Posons \(U=X-Y\), qui est une variété affine non sing.
\[
A_r^i(U)\simeq A_r^i(X)/\operatorname{Im}A_r^i(Y),\quad\text{de même pour }A_a^i(U)\ldots
\]
\[
\begin{array}{llll}
2i\leq n-2&\text{i.e. }i\leq\dfrac n2-1&A_i^r(U)\unsure{=|\equiv}0&\text{(exc. }p\text{ ?)}\deleted{\unsure{?}}\\
2i=n-1&\text{i.e. }i=\dfrac{n-1}{2}\deleted{\unsure{+1}}&A_i^a(U)=0&\text{(exc }p\text{ ?)}
\end{array}
\]

\textcircled{3'}\quad Soit \(E\) le cône projectant affine épointé de \(X\), (variété quasi-affine non sing. de dim \(2n+1\)).
\[
\begin{array}{lll}
2i\geq n+2&\text{i.e. }i\geq\dfrac n2+1&A_r^i(E)\ \deleted{\unsure{est}}\text{ fini},\quad A_r^i(E)^a=0\\
&n-i\leq\dfrac n2-1&\\[4pt]
2i\geq n+1&\text{i.e. }i\geq\dfrac{n+1}{2}&A_a^i(E)\ \text{fini}\\
&\text{i.e. }n-i\leq\dfrac{n-1}{2}&
\end{array}
\]

\uline{N.B.} On doit avoir
\[
\left\{\begin{array}{lll}
A_a^i(E)\simeq\deleted{\unsure{B_a^i(X)}}\inserted{\unsure{M_a^i(X)|N_a^i(X)}}&\inserted{\text{\tikz[baseline=(n.base)]\node[draw,ellipse,inner sep=1pt](n){mod groupes finis};}}&\text{si }2i\leq n\\
\deleted{A_{\unsure{r|\tau}}^i(E)^a\simeq B_\tau^i(X)}&\deleted{\text{mod grp finis}}&\deleted{\text{si }2i\leq n\unsure{+1}}\\
A_a^i(E)\simeq0&\text{mod grp finis}&\text{si }2i\geq n+1
\end{array}\right.
\]
\[
\left\{\begin{array}{lll}
A_r^i(E)^a\simeq\inserted{I^i}\deleted{\unsure{B_\tau^i}}(X)&\text{mod grp finis}&\text{si }2i\leq n+1\\
A_r^i(E)^{\unsure{a}}\simeq\unsure{0}&\text{``}&\text{si }2i\geq n+2
\end{array}\right.
\]

où
\[
\unsure{M_a^i|N_a^i}(X)=\operatorname{Ker}\bigl(A_a^i(X)\xrightarrow{L^{n-2i}}A_{\unsure{a}}^{n-i}(X)\bigr)
\]
\margin{partie ``primitive'' de la ``\unsure{ch.|cl.} algébrique''}
\[
\unsure{I^i|I_r^i}(X)=\operatorname{Ker}\bigl(A_r^i(X)^a\xrightarrow{L^{n-2i+1}}A_r^{n-i+1}(X)^a
\]
\margin{partie ``primitive'' de la \unsure{VAb.|V.A.} intermédiaire}
\margin{112}
